\section{ISO 27789: Audit Trails for Electronic Health Records}
\label{sec:iso27789}

ISO 27789:2021 berjudul \textit{Health Informatics --- Audit Trails for Electronic Health Records} merupakan standar internasional yang diterbitkan oleh ISO/TC 215 (\textit{Health Informatics}) bekerja sama dengan CEN/TC 251, sebagai edisi kedua yang menggantikan ISO 27789:2013. Standar ini menetapkan kerangka kerja bersama untuk jejak audit pada rekam kesehatan elektronik (\textit{Electronic Health Record}, EHR), mencakup definisi peristiwa pemicu audit (\textit{trigger events}) dan data audit, dengan tujuan agar seluruh informasi kesehatan pribadi (\textit{personal health information}) dapat ditelusuri jejaknya secara lengkap lintas sistem informasi maupun lintas domain organisasi. 

Standar ini berlaku untuk sistem yang memproses informasi kesehatan pribadi dan mewajibkan pembuatan rekam audit yang aman setiap kali pengguna membaca, membuat, memperbarui, atau mengarsipkan informasi tersebut, di mana rekam audit tersebut minimal harus dapat mengidentifikasi pengguna secara unik, mengidentifikasi subjek layanan (\textit{subject of care}) secara unik, mengidentifikasi fungsi yang dilakukan, serta mencatat tanggal dan waktu kejadian. Penting dicatat bahwa lingkup standar ini secara ketat dibatasi pada pencatatan peristiwa (\textit{logging of events}); perubahan nilai data pada field EHR diasumsikan tercatat pada sistem basis data EHR itu sendiri, bukan pada log audit, sehingga log audit hanya berisi identitas dan tautan (\textit{links}) ke segmen rekam medis, bukan isi informasi kesehatan pribadi itu sendiri.

Dari sisi persyaratan, ISO 27789 menetapkan sejumlah kebutuhan etis dan formal yang harus dipenuhi oleh jejak audit, di antaranya adanya kebijakan akses (\textit{access policy}) yang selaras dengan ISO 27799:2016 klausul 9.1.1, identifikasi pengguna sistem informasi secara tidak ambigu, pencatatan peran (\textit{role}) pengguna saat melakukan suatu aksi terhadap data, serta pembuatan rekam audit yang aman sesuai ISO 27799:2016 klausul 12.4.1. Standar ini juga menjabarkan kegunaan data audit untuk tiga kepentingan utama, yaitu tata kelola dan pengawasan (misalnya mendeteksi akses tidak sah, mengevaluasi akses darurat, dan mendeteksi penyalahgunaan hak akses), pemenuhan hak subjek layanan untuk mengetahui siapa saja yang telah mengakses rekam kesehatannya, serta pemenuhan kebutuhan bukti dan retensi data yang dapat digunakan sebagai bukti dokumenter oleh penyedia layanan kesehatan.

Selanjutnya, ISO 27789 mendefinisikan dua jenis peristiwa pemicu (\textit{trigger events}) yang wajib menghasilkan rekam audit, yaitu peristiwa akses (\textit{access events}) dan peristiwa kueri (\textit{query events}) terhadap informasi kesehatan pribadi, sementara peristiwa lain seperti proses \textit{start/stop} aplikasi, autentikasi pengguna, atau akses ke log audit itu sendiri berada di luar cakupan standar. Struktur rekam audit yang dirinci mencakup identifikasi peristiwa pemicu (ID peristiwa, kode aksi, waktu kejadian, indikator hasil, dan kode tipe peristiwa), identifikasi pengguna (ID pengguna, nama pengguna, kode peran, dan tujuan penggunaan data), identifikasi titik akses jaringan, identifikasi sumber audit, serta identifikasi objek partisipan yang diakses. Struktur ini diselaraskan dengan format DICOM sebagai salah satu pembaruan utama pada edisi 2021. Pada bagian akhir, standar ini juga mengatur pengelolaan data audit secara aman, mencakup ketersediaan sistem audit, persyaratan retensi, serta perlindungan kerahasiaan dan integritas jejak audit itu sendiri, sehingga log audit tidak dapat dengan mudah diubah atau dihapus oleh pihak yang tidak berwenang.

Secara keseluruhan, ISO 27789 memberikan kerangka acuan yang cukup rinci dan bersifat teknis mengenai apa saja yang perlu dicatat dalam sebuah audit trail rekam kesehatan elektronik, sekaligus mengapa pencatatan tersebut penting dari sisi etika, akuntabilitas, dan kepatuhan hukum. Standar ini melengkapi ISO 27799 (yang membahas manajemen keamanan informasi kesehatan secara umum) dengan memberikan spesifikasi konkret pada level elemen data audit, sehingga menjadi rujukan yang relevan bagi organisasi kesehatan maupun pengembang sistem informasi kesehatan dalam merancang mekanisme pencatatan aktivitas pengguna yang konsisten dan dapat dipertukarkan lintas domain.

\section{NIST SP 800-53 Revisi 5: Keluarga Kontrol Audit and Accountability (AU)}
\label{sec:nist80053au}

NIST Special Publication 800-53 Revisi 5 merupakan katalog kontrol keamanan dan privasi yang diterbitkan oleh \textit{National Institute of Standards and Technology} (NIST) untuk digunakan oleh sistem informasi dan organisasi, khususnya di lingkungan pemerintahan federal Amerika Serikat, namun juga diadopsi secara luas oleh berbagai industri termasuk sektor kesehatan. Dari 20 keluarga kontrol yang terdapat dalam dokumen ini, keluarga \textit{Audit and Accountability} (AU) secara khusus berfokus pada bagaimana suatu organisasi merencanakan, menghasilkan, melindungi, meninjau, serta menyimpan rekam audit (\textit{audit records}) sebagai bentuk pertanggungjawaban atas aktivitas yang terjadi pada sistem informasi. Keluarga kontrol ini diawali dengan AU-1 (\textit{Policy and Procedures}), yang mewajibkan organisasi untuk mengembangkan, mendokumentasikan, dan menyebarluaskan kebijakan serta prosedur audit dan akuntabilitas, termasuk menunjuk pejabat yang bertanggung jawab atas pengelolaannya dan meninjau kebijakan tersebut secara berkala atau setelah terjadi insiden tertentu.

Inti dari keluarga kontrol AU terletak pada kontrol AU-2 hingga AU-7, yang secara berurutan mengatur siklus hidup data audit mulai dari pembuatan hingga pemanfaatannya. AU-2 (\textit{Event Logging}) mensyaratkan organisasi mengidentifikasi jenis-jenis peristiwa yang perlu dicatat serta memberikan alasan mengapa jenis peristiwa tersebut memadai untuk mendukung investigasi pasca-insiden. AU-3 (\textit{Content of Audit Records}) menetapkan bahwa setiap rekam audit harus memuat informasi mengenai jenis peristiwa, waktu kejadian, lokasi kejadian, sumber peristiwa, hasil peristiwa, serta identitas individu atau entitas yang terlibat. Selanjutnya, AU-4 (\textit{Audit Log Storage Capacity}) mengatur alokasi kapasitas penyimpanan log audit sesuai kebutuhan retensi, AU-5 (\textit{Response to Audit Logging Process Failures}) mengatur mekanisme peringatan dan tindakan lanjutan apabila proses pencatatan audit mengalami kegagalan, AU-6 (\textit{Audit Record Review, Analysis, and Reporting}) mensyaratkan peninjauan dan analisis rekam audit secara berkala untuk mendeteksi aktivitas yang tidak wajar, dan AU-7 (\textit{Audit Record Reduction and Report Generation}) menyediakan kemampuan untuk merangkum serta menghasilkan laporan dari data audit tanpa mengubah isi maupun urutan waktu aslinya.

Kontrol berikutnya, yaitu AU-8 hingga AU-12, lebih menitikberatkan pada aspek integritas, kerahasiaan, dan ketertelusuran data audit. AU-8 (\textit{Time Stamps}) mewajibkan penggunaan cap waktu yang konsisten, umumnya mengacu pada \textit{Coordinated Universal Time} (UTC), agar urutan kejadian dapat direkonstruksi secara akurat. AU-9 (\textit{Protection of Audit Information}) mengatur perlindungan informasi audit dan perangkat pencatatannya dari akses, modifikasi, maupun penghapusan yang tidak sah, termasuk melalui penyimpanan pada media \textit{write-once}, pemisahan lokasi penyimpanan, maupun penerapan mekanisme kriptografi. AU-10 (\textit{Non-repudiation}) memastikan adanya bukti yang tidak dapat disangkal bahwa suatu individu telah melakukan tindakan tertentu, misalnya melalui tanda tangan digital, sementara AU-11 (\textit{Audit Record Retention}) mengatur jangka waktu penyimpanan rekam audit sesuai kebutuhan investigasi dan regulasi, dan AU-12 (\textit{Audit Record Generation}) memastikan kemampuan sistem untuk menghasilkan rekam audit pada komponen-komponen yang telah ditentukan secara konsisten dengan jenis peristiwa yang telah didefinisikan pada AU-2.

Adapun kontrol lanjutan dalam keluarga ini, yaitu AU-13 hingga AU-16, mencakup aspek yang lebih spesifik seperti pemantauan terhadap pengungkapan informasi (\textit{Monitoring for Information Disclosure}), audit terhadap sesi pengguna (\textit{Session Audit}) yang memungkinkan perekaman aktivitas seorang pengguna secara mendetail dalam kondisi tertentu, serta audit lintas organisasi (\textit{Cross-organizational Audit Logging}) yang menjadi relevan ketika data audit perlu dikoordinasikan antarorganisasi yang berbagi sistem atau layanan. Secara umum, keluarga kontrol AU pada NIST SP 800-53 Revisi 5 ini bersifat lebih generik dan berorientasi manajerial-teknis dibandingkan ISO 27789, karena dirancang untuk berlaku lintas domain sistem informasi apa pun, tidak terbatas pada rekam kesehatan elektronik. Meski demikian, prinsip-prinsip yang ditekankan seperti identifikasi pengguna, pencatatan waktu kejadian, perlindungan integritas log, serta retensi data audit memiliki kesesuaian yang tinggi dengan persyaratan pada ISO 27789, sehingga kedua standar ini dapat saling melengkapi sebagai acuan dalam merancang mekanisme audit trail pada sistem informasi kesehatan.

\section{NIST SP 800-92 --- Guide to Computer Security Log Management}

NIST Special Publication (SP) 800-92, \textit{Guide to Computer Security Log Management}, diterbitkan oleh National Institute of Standards and Technology sebagai panduan teknis operasional untuk pengelolaan log keamanan komputer. Berbeda dengan ISO 27789 yang berfokus pada persyaratan audit trail spesifik rekam medis elektronik dan NIST SP 800-53 Rev. 5 yang menetapkan kontrol pada tataran kebijakan dan tata kelola (governance), SP 800-92 beroperasi pada tataran implementasi teknis: bagaimana log dihasilkan, ditransmisikan, disimpan, dianalisis, dan dihapus secara efektif dalam suatu \textit{log management infrastructure}. Dokumen ini mendefinisikan log sebagai catatan atas peristiwa yang terjadi dalam sistem atau jaringan organisasi, serta menekankan bahwa pengelolaan log yang baik harus menyeimbangkan antara sumber daya pengelolaan yang terbatas dengan volume data log yang terus bertambah. Dalam konteks penelitian ini, SP 800-92 digunakan sebagai acuan teknis untuk menilai kelayakan operasional mekanisme audit trail pada sistem DecMed, melengkapi kerangka kepatuhan yang telah dibahas pada dua subbab sebelumnya.

SP 800-92 menguraikan siklus hidup pengelolaan log ke dalam lima tahapan utama, yaitu \textit{generation}, \textit{transmission}, \textit{storage}, \textit{analysis}, dan \textit{disposal}. Pada tahap penyimpanan, dokumen ini menekankan bahwa integritas dan ketersediaan log harus dijaga dari perubahan maupun penghapusan yang tidak sah, antara lain melalui pembatasan hak akses (\textit{append-only privileges}), penggunaan \textit{message digest} atau fungsi hash untuk memverifikasi bahwa berkas log tidak mengalami modifikasi selama penyimpanan maupun transfer, serta enkripsi untuk melindungi kerahasiaan data log yang mengandung informasi sensitif. Ketiga aspek ini---kerahasiaan, integritas, dan ketersediaan (\textit{confidentiality, integrity, availability})---menjadi kriteria evaluasi yang relevan untuk menilai apakah suatu arsitektur penyimpanan log, termasuk arsitektur penyimpanan terdesentralisasi, memenuhi prinsip dasar pengelolaan log yang aman.

Terkait retensi dan pengarsipan, SP 800-92 tidak menetapkan durasi retensi tunggal yang bersifat universal, melainkan mengarahkan organisasi untuk menentukan periode retensi berdasarkan kombinasi antara persyaratan regulasi yang berlaku, kebutuhan operasional, dan hasil analisis risiko masing-masing organisasi. Dokumen ini juga mengingatkan bahwa pemilihan media dan format penyimpanan arsip harus mempertimbangkan daya tahan media terhadap periode retensi yang ditetapkan, serta perlunya peninjauan berkala untuk memastikan data arsip tetap dapat diakses dan tidak berisiko usang secara teknologi. Selaras dengan arahan tersebut, penelitian ini menetapkan durasi retensi 5 (lima) tahun sebagai baseline asumsi awal untuk keperluan evaluasi, dengan ketentuan bahwa pada tahap implementasi durasi retensi tersebut bersifat dapat dikonfigurasi oleh administrator sistem sesuai kebutuhan operasional maupun regulasi yang berlaku, sejalan dengan prinsip fleksibilitas kebijakan retensi yang dianjurkan oleh SP 800-92.

Kerangka SP 800-92 relevan digunakan untuk menilai kelayakan infrastruktur penyimpanan DecMed dari sisi operasional pengelolaan log, di luar aspek performa penyimpanan yang berada di luar batasan penelitian ini. Arsitektur DecMed yang membagi data menjadi komponen \textit{on-chain} pada IOTA Tangle untuk data administratif dan \textit{capability entries} yang bersifat \textit{immutable}, komponen \textit{off-chain} pada IPFS Cluster untuk data klinis terenkripsi berukuran besar sehubungan dengan batas ukuran objek IOTA sebesar 250 KB, serta Redis untuk data sementara pada proses \textit{proxy re-encryption} (PRE), dapat dipetakan terhadap prinsip pengelolaan log SP 800-92. Sifat \textit{immutable} pada Tangle secara konseptual sejalan dengan tujuan verifikasi integritas melalui \textit{message digest} yang disyaratkan SP 800-92, sementara pemisahan data administratif dan data klinis pada dua lapisan penyimpanan yang berbeda perlu dikaji lebih lanjut untuk memastikan bahwa mekanisme kerahasiaan, integritas, dan ketersediaan log audit tetap terjaga secara konsisten pada kedua lapisan tersebut, termasuk kejelasan kebijakan retensi dan penghapusan data pada IPFS Cluster maupun Redis. Kajian lebih mendalam mengenai kesesuaian arsitektur ini terhadap ketiga kerangka acuan---ISO 27789, NIST SP 800-53 (AU), dan NIST SP 800-92---akan disintesis pada bagian analisis penelitian ini.